\section{现代开放对齐生态：更多玩家，仍然缺数据}

前几章按时间追踪模型，本章把 2024 年 4 月视为课堂快照。参与者从大学、非营利机构、Hugging Face 社区扩展到大量公司；模型 recipe 更复杂，评测更快，开放/闭源差距在部分任务缩小。但讲者认为最稳定的限制仍是数据：尤其是合法、开放、多样的人类 preference data。

\subsection{更多模型、更多数据生成角色}

本节展示现代生态的多样性。不同团队专门做 instruction rewriting、preference generation、过滤、chat data、reward model 和训练框架。开放 alignment 不再是单团队端到端工作，而是 artifact supply chain。

\lecturefigure{slide-071.jpg}{现代生态中的数据生成、模型训练与社区角色多样化}{official deck, p.~71}

\begin{knowledgebox}{Artifact supply chain}
一个开放 aligned model 可能依赖：第三方 base weights、另一个模型生成的 instructions、公开 chat logs、社区 preference data、开源 trainer、闭源 LLM judge 与公开 leaderboard。每个依赖都会引入许可、污染、风格和可复现性风险，因此 model card 应记录完整 lineage。
\end{knowledgebox}

\subsection{Llama 3：更多是 scale story，也是 post-training signal}

Llama 3 在课堂前后发布，讲者把它视为“更多关于 scaling”而非单一 alignment breakthrough。Meta 报告使用 IFT、rejection sampling、DPO 和 PPO 等多阶段方法；复杂 pipeline 说明行业愿意叠加方法获得小幅收益，也说明尚未找到统一简单算法。

\lecturefigure{slide-072.jpg}{Llama 3：大规模 pretraining 与多阶段 post-training 的课堂快照}{official deck, p.~72}

Rejection sampling 可概念性写成：从 policy 采样 $K$ 个候选
\[
y_1,\ldots,y_K\sim\pi(\cdot\mid x),
\qquad y^{\star}=\arg\max_{y_i} r(x,y_i),
\]
其中 $r$ 可以是 reward model、规则或 judge。把 $y^{\star}$ 回灌 IFT 能提升初始化，但也会把 reward/judge 偏差固化进新数据。

\teachervoice{01:02:44--01:03:22，Nathan 对 Llama 3 同时使用 IFT、rejection sampling、DPO 与 PPO 感到配方复杂；他猜测这些阶段在逐步移动能力并为下一阶段初始化，未来可能被蒸馏成更简单算法。}

\subsection{Current directions：alignment 正在变成 data problem}

本节从模型转向研究问题。开放社区只有少量广泛使用的 preference datasets，例如 Anthropic HH、UltraFeedback 和 Nectar；重复使用会导致过拟合、同质化和 contamination。未来需要更多真实人类、多语言、多领域和长期交互数据。

\lecturefigure{slide-073.jpg}{Current directions：开放语言模型对齐的研究前沿}{official deck, p.~73}

\lecturefigure{slide-074.jpg}{开放与闭源 aligned models 的能力差距随任务而变化}{official deck, p.~74}

\begin{warningbox}{“Open catches closed” 必须指定任务和日期}
某些开放模型在公开 benchmark 上接近闭源模型，不代表工具使用、长上下文、可靠性、安全、延迟和多语言都追平。图是 2024 年 4 月 snapshot；模型版本和 judge 更新后结论会改变。比较必须给出 exact model、date 与 evaluation protocol。
\end{warningbox}

\lecturefigure{slide-075.jpg}{当前方向一：Preference data 严重不足，数据质量限制实验}{official deck, p.~75}

\begin{importantbox}{Preference dataset 的覆盖矩阵}
高质量数据至少应覆盖：任务域、语言、用户群体、风险级别、回答长度、工具/无工具、多轮状态和 disagreement。只扩大同一 synthetic template 的样本数，会降低方差却不扩大行为边界。
\end{importantbox}

\teachervoice{01:00:00--01:03:48，讲者把 preference data scarcity 作为当前核心瓶颈，并提醒 synthetic data 常重复、分布不稳，因而对 alignment 的泛化提升有限。}

\subsection{开放对齐发生在哪里}

最后一张内容页列出 AI2、HuggingFaceH4、Berkeley、LMSYS、NousResearch 等参与者。名单的意义不是完整排名，而是说明开放进展分散在模型、数据、trainer、evaluation 与社区运营多个节点；健康生态需要这些 artifact 能互相组合并被审计。

\lecturefigure{slide-076.jpg}{开放 alignment 的主要社区、机构与 artifact 类型}{official deck, p.~76}

\begin{knowledgebox}{读图：机构名单背后的分工}
AI2 强调 Tulu/OLMo 与数据发布，HuggingFaceH4 快速复现 recipe，LMSYS 提供 Arena 与真实对话数据，Berkeley/社区团队探索 PPO 和评测，其他团队生成 instructions 或 preference data。开放生态的优势是并行探索，弱点是 provenance、许可和评测口径不统一。
\end{knowledgebox}

\teachervoice{01:04:49--01:05:30 的研究建议是：没人能跟踪所有进展，应该持续构建技能和自己认为重要的东西。对开放生态而言，这意味着贡献高质量数据、评测或工具同样是核心研究，不必只追模型规模。}

\subsection{Synthetic data 与人类数据的关系}

合成数据能绕过部分版权、隐私和标注成本，也能生成更大规模 token；但 teacher model 的偏差、重复模式和缺少真实用户 intent 会限制泛化。更合理的方向是人类 seed、模型扩增、自动过滤和人工审计的混合 pipeline。

设真实数据分布为 $p_h$、合成分布为 $p_s$，混合训练分布可写成
\[
p_{\lambda}=\lambda p_h+(1-\lambda)p_s,
\]
其中 $\lambda$ 控制人类数据比例。最优 $\lambda$ 不是固定常数：数据质量、任务域和 teacher strength 改变时，需要通过 held-out human evaluation 重新选择。

\teachervoice{Q\&A 的 01:14:38--01:15:29，Nathan 认为 synthetic data 是少数能继续生成更多 token 的路线，也相信模型会持续参与训练其他模型；但他明确说这一方向仍很早。}

\subsection{本章小结}

现代生态的主角从单一模型扩展成 artifact network。Llama 3 展示大规模 base 与多阶段 post-training 的组合；开放/闭源差距依任务变化；高质量 preference data 和可信评测仍是稀缺资源。未来进展可能来自更简单的统一算法，也可能来自更好的数据与反馈基础设施。

